%% Abstract A – Forward surrogate (pdfLaTeX compatible)
%% Overleaf default compiler (pdfLaTeX) works without any changes.
\documentclass[11pt,a4paper]{article}

\usepackage[a4paper, top=1in, bottom=1in, left=1in, right=1in]{geometry}
\usepackage[utf8]{inputenc}
\usepackage[T1]{fontenc}
\usepackage[english]{babel}
\usepackage{newtxtext,newtxmath}   % Times-style body + math fonts
\usepackage{microtype}
\usepackage{setspace}
\setstretch{1.12}
\usepackage{booktabs}
\usepackage[hidelinks]{hyperref}

\setlength{\parindent}{0pt}
\setlength{\parskip}{0.55em}
\widowpenalty=10000
\clubpenalty=10000

\usepackage{titling}
\setlength{\droptitle}{-1.8em}
\pretitle{\begin{center}\large\bfseries}
\posttitle{\par\vskip0.2em\hrule height 0.4pt\par\end{center}\vskip1em}
\preauthor{}\postauthor{}
\predate{}\postdate{}

\begin{document}

\title{Spectrally-consistent neural surrogate for real-time ultrasound
       pressure-field prediction in heterogeneous breast tissue}
\date{}
\maketitle

\noindent
\textbf{Background.}
Patient-specific HIFU planning requires fast pressure-field simulation in
heterogeneous tissue.
Full-wave k-Wave delivers reference accuracy but costs
\textbf{{\raise0.5ex\hbox{$\sim$}}60\,s per 316$\,\times\,$256\,/\,1\,MHz
configuration} on an RTX\,4070---an order of magnitude too slow for the
sub-second feedback clinicians need while sweeping the focal target.
Frequency-domain Helmholtz, ray tracing and lookup tables each fail on a
different axis (transient information, refraction tail, or per-patient
adaptivity), defining the engineering slot for a learned forward surrogate that
is on-the-go and GPU-light.

\noindent
\textbf{Method.}
We trained a 2-D Fourier Neural Operator (FNO) and a matched 2-D U-Net
baseline on \textbf{1\,000 simulations} drawn from the recently released
\textbf{OpenBreastUS} dataset (Zeng \textit{et al.}\ 2025)---anatomically
realistic sound-speed, density and attenuation maps derived from real breast
MRI, chosen for \textit{in\,vivo} fat\,/\,fibroglandular impedance contrast that
synthetic phantoms cannot reproduce.
Identical 70\,/\,15\,/\,15 split (seed\,0) across all backbones; combined
\texttt{LpLoss\,+\,0.3$\cdot$H1Loss} objective; per-channel $z$-scored inputs;
\texttt{log1p(}$p_{\max}$\texttt{)} $z$-scored targets recovered to Pa via
\texttt{expm1}.

\noindent
\textbf{Results.}
FNO achieves \textbf{test relative L2\,=\,0.097}, against U-Net's
\textbf{0.264} (\textbf{2.7$\times$} lower error) and a ConvNeXt-2D
ablation's \textbf{0.99} (an order of magnitude worse).
Inference completes in \textbf{{\raise0.5ex\hbox{$\sim$}}8\,ms\,/\,sample}
on the same RTX-4070---a \textbf{{\raise0.5ex\hbox{$\sim$}}7\,500$\times$
speed-up} over k-Wave at matched accuracy on the test split.
Validation loss scaled $0.81\to0.39$ from $200\to1{,}000$ samples, indicating
substantial headroom on the full 8\,000-phantom set.
Visual inspection confirms FNO recovers the refraction tail that U-Net smears,
consistent with FFT-based representations matching wave-equation structure.

\noindent
\textbf{Significance.}
To our knowledge this is the first application of OpenBreastUS to ultrasound
treatment-planning surrogates (the original release targets ultrasound CT
reconstruction).
Prior neural-operator work (TUSNet 2022, DeepTFUS 2023, Stanziola 2023) used
non-breast tissue or synthetic phantoms; the \textbf{spectrally-consistent
backbone advantage} (FNO vs.\ generic CNNs at matched parameter budgets) is a
heterogeneous-tissue-specific finding not previously established and a
clinically relevant building block for real-time HIFU planning.

\smallskip\noindent
\textit{Keywords:} HIFU, Fourier Neural Operator, k-Wave, OpenBreastUS,
real-time treatment planning, wave equation, breast ultrasound.

\end{document}
